% ---------------- Title page ----------------
\begin{titlepage}
\thispagestyle{empty}
\centering\setstretch{1.2}
{\Large\bfseries Project Report On\par}
\vspace{14pt}
{\LARGE\bfseries \thesistitle\par}
\vspace{8pt}
{\large A Cross-Sectional Study\par}
\vspace{22pt}
% TODO: Add or confirm the university logo as figures/logo.png.
\IfFileExists{figures/logo.png}{\includegraphics[width=1.7in]{figures/logo.png}\par}{{\fbox{\parbox[c][1.3in][c]{1.7in}{\centering\small\color{red}\textbf{University logo:}\\ place logo.png\\ in figures/}}}\par}
\vspace{22pt}
A Dissertation Submitted to the \departmentname, \institution, for the Partial Fulfilment of the Requirements for the Degree of \degreename\par
\vspace{26pt}
{\large\bfseries Submitted By\par}
\vspace{6pt}
{\bfseries \authorname\par}
ID: \authorid\par
Batch: \authorbatch\quad Section: \authorsection\par
Program: \degreename\par
\departmentname\par
\institution\par
\institutionaddress\par
\vspace{18pt}
{\large\bfseries Supervised By\par}
\vspace{6pt}
\supervisorname\par
\supervisortitle\par
\institution\par
\vfill
\end{titlepage}

\clearpage
\pagenumbering{roman}
\setcounter{page}{1}

% ---------------- Certification ----------------
% ==========================================
% PAGE 4: CERTIFICATION
% ==========================================
\begin{center}
    \textbf{\large CERTIFICATION}
\end{center}
\vspace{1em}

This is to certify that the research work embodied in this project entitled ``\textbf{Acceptance of HPV Vaccination Among Female Students}'' is the original and independent work carried out by \textbf{Sadia Benta Alam Bushra (2022200300048)} under our supervision.

This project has not been submitted to any other university or institution for the award of any degree/diploma to the best of our knowledge. All sources of information used in this project have been properly acknowledged and referenced within the text.

This is to certify that this thesis is submitted as a partial requirement for the Bachelor of Pharmacy (B. Pharm) degree at the Department of Pharmacy, Southeast University, Dhaka, Bangladesh.

\vspace{8cm}

\noindent
\textbf{Professor Dr. Hasnat M Alamgir, PhD}\\
Director, IRT\\
Department of Pharmacy\\
Southeast University\\
251/A Tejgaon I/A, Dhaka 1208, Bangladesh

\vfill
\noindent\rule{\textwidth}{0.5pt}
\clearpage

% ---------------- Declaration ----------------
\chapter*{Declaration}
\addcontentsline{toc}{chapter}{Declaration}
I do here by declare that the research work embodied in this project entitled ``\textit{Acceptance of HPV Vaccination Among Female
Students }'', submitted to Southeast University, Dhaka, Bangladesh in partial fulfillment of the requirements for the Degree of Bachelor of Pharmacy, is my original work.

I further declare that this study has not been previously submitted to any other university or institution for the award of any degree or diploma.

\vspace{8cm}

\noindent
\textbf{SADIA BENTA ALAM BUSHRA}\\
\textbf{ID: 2022200300048}\\
\textbf{Batch : 40}\\
\textbf{Section: B}\\
\textbf{Student, Department of Pharmacy (Hons.)}\\
\textbf{Southeast University}\\
\textbf{251/A Tejgaon I/A, Dhaka 1208, Bangladesh}

\vfill
\noindent\rule{\textwidth}{0.5pt}
\clearpage

% ---------------- Letter of transmittal ----------------
\chapter*{Letter of Transmittal}
\addcontentsline{toc}{chapter}{Letter of Transmittal}
\submissionmonth\par
\supervisorname\par
\supervisortitle\par
\institution\par
\institutionaddress\par

\textbf{Subject: Submission of project report on ``\thesistitle''.}

Dear Sir,

I am pleased to submit my project report entitled ``\thesistitle: A Cross-Sectional Study'' in partial fulfilment of the requirements for the degree of \degreename. The report presents the findings of a survey of female university students on their knowledge of, attitudes towards, and acceptance of HPV vaccination.

I am grateful for your guidance and support throughout this work. I would be happy to answer any questions about the report.

Sincerely,

\vspace{30pt}
\noindent
\begin{minipage}[t]{0.48\textwidth}
\end{minipage}\hfill
\begin{minipage}[t]{0.48\textwidth}
\centering
\rule{6cm}{0.4pt}\\
\textbf{\authorname}\\
ID: \authorid\\
\departmentname, \institution
\end{minipage}
\clearpage

% ---------------- Acknowledgment ----------------
\chapter*{Acknowledgment}
\addcontentsline{toc}{chapter}{Acknowledgment}
Above all, I extend my profound gratitude to Almighty Allah, whose infinite blessings, guidance, and mercy bestowed upon me the perseverance, capability, and fortitude to overcome every obstacle and complete this research project successfully.

I sincerely thank my respected supervisor, \textbf{Dr. Hasnat M Alamgir, PhD, Department of Pharmacy, Director of IRT, Southeast University}, and my co-supervisor \textbf{Dr. Sabiha Ferdowsy Koly Assistant Professor, Southeast University} for theme invaluable guidance, constructive feedback, and continuous encouragement throughout this research. His suggestions and support helped me better understand the research process and complete this project.

I am also deeply grateful to \textbf{Dr. Israt Jahan Bulbul}, Chairman and Associate Professor of the Department of Pharmacy at Southeast University, for her constant encouragement and valuable support. My sincere thanks also go to \textbf{Dr. Nusratun Nahar}, Associate Professor and REB Board Member, for her crucial assistance in granting ethical approval for this study. Additionally, I truly appreciate all the faculty members of the department for their continuous guidance throughout my academic journey.

I would like to extend my heartfelt thanks to our senior, \textbf{Afrin Nur}, for her constant guidance and immense help from the very beginning of this project. My deep appreciation also goes to my incredible batchmate \textbf{MD. Kajul Mia} for his constant support and coordination. Last but not least, I am profoundly grateful to all the participants of this study. Without their willingness to share their experiences, this research would not have been possible.

I would also like to thank my \textbf{friends and classmates, faculty members, and staff of the Department of Pharmacy, Southeast University}, for their encouragement, cooperation, and moral support during my academic journey.

Finally, I acknowledge everyone who directly or indirectly contributed to the successful completion of this research. Their support and encouragement have been truly invaluable to me.

\vfill
\noindent\rule{\textwidth}{0.5pt}
\clearpage

% ---------------- Abstract ----------------
\chapter*{Abstract}
\addcontentsline{toc}{chapter}{Abstract}
\onehalfspacing
\textbf{Background:} Human papillomavirus (HPV) causes almost all cervical cancers, which are among the most common cancers in women in Bangladesh. The national HPV vaccination campaign has so far targeted girls aged 10--14 years, leaving university-age women outside the programme. Little is known about whether they would accept the vaccine. This study measured acceptance of HPV vaccination among female university students in Dhaka and identified the factors associated with it.

\textbf{Methods:} A cross-sectional survey was conducted among 231 female students aged 18--28 years from universities in Dhaka (September--October 2026), using a structured questionnaire. Acceptance was defined as having received at least one HPV vaccine dose or being ``likely'' or ``very likely'' to be vaccinated in future. Knowledge (15 items, score 0--15; KR-20 = 0.828) and attitude (10 items, score 10--50; Cronbach's $\alpha$ = 0.823) were also scored. Associations with acceptance were tested with chi-square, Fisher's exact and Mann--Whitney $U$ tests, and multivariable logistic regression was used to estimate adjusted odds ratios (AOR).

\textbf{Results:} Of 229 students with a classifiable outcome, 184 (80.3\%) were accepting: 39 (17.0\%) had already received at least one dose, and 145 of 190 unvaccinated or unsure students (76.3\%) were likely or very likely to be vaccinated. Acceptance was almost identical in healthcare (79.8\%) and non-healthcare (80.7\%) students. Although 79.8\% had heard of HPV, knowledge was low (mean 6.32 $\pm$ 3.67 out of 15; 67.1\% poor) while attitudes were highly positive (mean 44.69 $\pm$ 4.48 out of 50; 88.3\% positive). The most common reason for not being vaccinated was not knowing about the vaccine (52.3\%). In adjusted analysis, having heard of HPV (AOR 3.31, 95\% CI 1.44--7.61) and a higher attitude score (AOR 1.08 per point, 95\% CI 1.00--1.17) were associated with acceptance. Free or subsidised vaccine (45.9\%), a doctor's recommendation (42.4\%) and a university education campaign (38.5\%) were the factors most often named as encouraging vaccination.

\textbf{Conclusion:} Most female university students in this sample would accept HPV vaccination, but few had been vaccinated, mainly because they did not know about the vaccine. Awareness campaigns on campus, provider recommendation and affordable access are likely to convert intention into uptake. Because the sample was non-random and only 45 students were classed as not accepting, the associations should be treated as exploratory.

\textbf{Keywords:} human papillomavirus; HPV vaccine; vaccine acceptance; female university students; Bangladesh
\clearpage

% ---------------- Contents ----------------
\setstretch{1.15}
\tableofcontents
\clearpage\listoftables\addcontentsline{toc}{chapter}{List of Tables}
\clearpage\listoffigures\addcontentsline{toc}{chapter}{List of Figures}
